\subsection{Optimization horizon and the economic meaning of a move}
The distinction between a one-period objective and a policy's holding value
matters even with known moments. The following scalar examples specialize
standard cost-aware trading calculations \cite{boyd} to the comparisons being
tested; they are not new solutions to the general dynamic trading problem.
There is no terminal liquidation charge in either example, matching the
simulation's horizon accounting.

\begin{proposition}[Parking can overcome myopic inertia]
Let returns be iid with mean zero, variance $\sigma^2>0$ and a finite fourth
moment. An incumbent
allocation $0<b\le L$ satisfies $\gamma\sigma^2b\le c$. With known moments,
the one-period optimizer on the full scalar support keeps $b$ forever.
Comparing immediate parking at zero with that continuing policy, the expected
sum of utility differences over $H$ periods is
\begin{equation}
 H\frac{\gamma\sigma^2b^2}{2}-cb-\frac\eta2b^2.
 \label{eq:park-horizon}
\end{equation}
The first-period gain is negative, while the holding-period gain is positive
if $H>2c/(\gamma\sigma^2b)+\eta/(\gamma\sigma^2)$.
\end{proposition}
\begin{proof}
At the unchanged target, the smooth marginal is $-\gamma\sigma^2b$, within
the proportional-cost subgradient band. Concavity therefore makes $b$ optimal
for the one-period problem. Repeating that rule preserves its state. Parking
pays $cb+\eta b^2/2$ once and then has zero utility; continuing at $b$ has no
turnover cost and mean utility $-\gamma\sigma^2b^2/2$ each period. Subtraction
gives \eqref{eq:park-horizon}. Its first-period value is at most
$-cb/2-\eta b^2/2<0$ under the stated condition.
\end{proof}

For $\gamma=5$, $c=\eta=10^{-4}$, $\sigma=.002$, $b=.5$ and $H=252$,
the deployment-day gain is $-.600$ basis points of utility, while the expected
holding-period gain is $5.675$ basis points. Thus a negative deployment-day
diagnostic is compatible with useful parking. This does not contradict support
monotonicity: the comparison is with a repeatedly myopic policy. A dynamically
optimized full-support policy could itself choose to liquidate. The two shadow
states also differ after parking, exactly the distinction retained in
\eqref{eq:bridge}. Stationarity here supplies the future holding calculation;
an episode rejection alone does not supply it.
The continuing rule has zero state-matched one-period oracle regret at every
date. Parking has positive regret $cb+\eta b^2/2-\gamma\sigma^2b^2/2$ on
deployment and zero thereafter, yet has higher expected cumulative utility
when \eqref{eq:park-horizon} is positive. State-matched regret measures local
implementation efficiency; it cannot replace full-path performance or turn
the one-period Oracle into an upper bound on cumulative policy value.

\begin{proposition}[A profitable horizon trade can be absent from the tested pair]
Let iid returns with a finite fourth moment have $0<\mu\le c$ and raw second
moment $M>0$. Starting at
zero, the known-moment one-period optimizer remains at zero, so adding the
scalar support leaves both comparison policies identically in cash. Yet, if
$H\mu>c$, a positive constant position held for $H$ periods can have positive
expected cumulative utility after its initial entry cost. Within the class of
constant nonnegative positions, its optimum is
\begin{equation}
 b_H=\min\left\{L,\frac{(H\mu-c)_+}{H\gamma M+\eta}\right\}>0.
 \label{eq:horizon-entry}
\end{equation}
\end{proposition}
\begin{proof}
The one-period gradient at zero is $\mu\le c$, so the proportional-cost band
prevents entry and the state remains zero. A constant position has expected
cumulative utility $(H\mu-c)b-(H\gamma M+\eta)b^2/2$. Maximizing this concave
quadratic on $[0,L]$ yields \eqref{eq:horizon-entry}; its positive right
derivative at zero ensures a positive value when $H\mu>c$.
\end{proof}

With $\mu=.00005$, return standard deviation $.01$ and the same $\gamma,c,
\eta,H$ as above, $M=.0001000025$, $b_H=.099125$ and the constant-position
value is $6.195$ basis points of cumulative utility. The one-period comparison
still has $Y=0$. The horizon trade belongs to a different policy rule; valid
inference for this zero-score pair does not certify that other rule's value.
Horizon-aware candidate construction is compatible with T1 when
it is predictable and costs are correctly accounted for, but it would change
the financial experiment and is not silently substituted in the reported
one-period design.
